Deux sons de même amplitude, de fréquences \qty{\v{f1}}{Hz} et \qty{\fpeval{\v{f1}+\v{df}}}{Hz}, se superposent en un point. Le déplacement résultant peut s'écrire
\[ y = \left[2A\cos\left(2\pi\,\frac{f_2 - f_1}{2}\,t\right)\right]\sin\left(2\pi\,\frac{f_1 + f_2}{2}\,t\right) \]

\begin{enumerate}[label=\alph*)]
    \item Quelle est la fréquence du son perçu ?
    \item Combien de temps s'écoule entre deux maximums d'intensité consécutifs ? Expliquez pourquoi la fréquence des battements vaut $f_2 - f_1$ et non $(f_2 - f_1)/2$.
    \item Combien de battements entend-on en \qty{\v{t}}{s} ?
\end{enumerate}
